% Related Work - ICML style

\paragraph{Trajectory prediction versus route generation.}
Short-horizon trajectory prediction forecasts positions seconds to minutes ahead and evaluates point-wise displacement error \citep{gupta2018social,salzmann2020trajectronpp,alahi2016social,zhou2022hivt}.
At these time scales, corridor commitment is rarely necessary---the vehicle remains within a single corridor throughout the prediction window.
Route generation operates at the trip level (minutes to hours), where the central challenge shifts from ``where will the vehicle be in five seconds'' to ``which of several topologically distinct corridors will the driver choose.''
This distinction is critical: metrics designed for short-horizon accuracy (ADE, FDE) do not capture corridor-level multi-modality.

\paragraph{Sequence-level route generation.}
Autoregressive models that predict the next node or GPS point are a natural formulation for route generation \citep{kong2023mobility_survey,hsu2024trajgpt}.
Diffusion-based approaches generate trajectories in continuous coordinate space \citep{zhu2023difftraj,wei2024diffrntraj,zhu2024controltraj,yang2025geo}, sometimes incorporating road-network constraints through topology-aware denoising.
These methods share a common limitation: corridor choice is represented only implicitly---either buried in the hidden state of an AR model or distributed across the full denoising trajectory.
When corridor alternatives are far apart in output space, implicit representation leads to mode averaging or long-horizon drift.

\paragraph{Latent variable models for trajectories.}
VAE-based and flow-based generative models capture multi-modality through latent variables \citep{kingma2014auto}.
Foundation models for mobility aim to learn reusable trajectory representations \citep{lin2024trajfm,unitraj2025,choudhury2024trajectorypowered,liu2025trajmamba}.
However, none of these approaches address the capacity matching between the latent space and the target modality structure: if the latent dimension is too high relative to the number of meaningful modes, the generative model must still navigate empty inter-modal regions, recovering the original mode-averaging problem.
CascadeTraj identifies this capacity matching as a necessary condition for tractable conditional generation.

\paragraph{Coarse-to-fine and hierarchical generation.}
Hierarchical strategies decompose generation into coarse structure and fine detail.
Cardiff \citep{guo2025cardiff} uses cascaded hybrid diffusion at multiple spatial resolutions; DiffPath \citep{fan2024diffpath} applies latent diffusion for road-network paths; SEED \citep{rao2025seed} bridges sequence and diffusion models; GTG \citep{wang2025gtg} demonstrates cross-city trajectory generation.
CascadeTraj shares the coarse-to-fine intuition but differs in a key respect: the coarse stage is an \emph{explicit corridor commitment} in a capacity-matched latent space, not a spatial downsampling or a latent code without controlled information content.
The bottleneck dimensionality is chosen to retain corridor identity while discarding intra-corridor variation, making the flow model's conditional generation task tractable by design.
